\paragraph{From decoded cochains to fields on a solid.}
The twelve vertices, thirty edges and twenty faces form a closed surface.
A concrete volume reconstruction is obtained by coning its declared
Euclidean realization to the centre. Use the registered coordinates
\((0,\pm1,\pm\varphi)\), \((\pm\varphi,0,\pm1)\),
\((\pm1,\pm\varphi,0)\), with their fixed port labels and
\(\varphi=(1+\sqrt5)/2\). Each boundary edge has length two; each of the
twenty positively oriented tetrahedra has volume \((3+\sqrt5)/6\).
The cone has \(13,42,50,20\) simplices in degrees zero through three.
This geometric realization and its numerical length unit are supplied.

\begin{theorem}[Gauge-covariant cone reconstruction]
\label{thm:cone-whitney-bridge}
Let \(D\) and \(C\) denote the signed boundary gradient and curl.
For \(\Pi_m=\mathbf1_m\mathbf1_m^{\mathsf T}/m\), set
\[
 G=(D^{\mathsf T}D+\Pi_{12})^{-1}-\Pi_{12},\qquad
 H=(CC^{\mathsf T}+\Pi_{20})^{-1}-\Pi_{20}.
\]
Order cone edges radially first and cone faces by boundary faces first.
Then its incidence maps and trace-preserving extensions are
\begin{align*}
 \widetilde D&=\begin{pmatrix}-\mathbf1_{12}&I_{12}\\0&D\end{pmatrix},&
 \widetilde C&=\begin{pmatrix}0&C\\-D&I_{30}\end{pmatrix},&
 \widetilde B&=\begin{pmatrix}I_{20}&-C\end{pmatrix},\\
 P_0&=\begin{pmatrix}\mathbf1_{12}^{\mathsf T}/12\\I_{12}\end{pmatrix},&
 P_1&=\begin{pmatrix}GD^{\mathsf T}\\I_{30}\end{pmatrix},&
 P_2&=\begin{pmatrix}I_{20}\\C^{\mathsf T}H\end{pmatrix}.
\end{align*}
They satisfy
\[
 \widetilde DP_0=P_1D,\qquad
 \widetilde CP_1=P_2C,\qquad
 \widetilde BP_2=\Pi_{20}.
\]
A boundary two-cochain has a closed cone extension if and only if its
total flux vanishes. In particular, every \(B=CA\) has such an extension.

Let \(W_k\) be Whitney interpolation on this geometric cone. On a uniform
time grid \(t_n=nh\), \(h>0\), interpolate \(P_1A_n\) linearly and
\(P_0\phi_n\) constantly on each open slab. For \(s=(t-t_n)/h\), the fields are
\begin{align*}
 \mathcal E(t)&=W_1P_1\left[-(A_{n+1}-A_n)/h-D\phi_n\right],\\
 \mathcal B(t)&=W_2P_2\left[(1-s)CA_n+sCA_{n+1}\right].
\end{align*}
The discrete gauge rule lifts to the ordinary potential gauge rule.
On the interior of the supplied solid and time window,
\(\partial_t\mathcal B+d\mathcal E=0\) and \(d\mathcal B=0\)
hold distributionally.
\end{theorem}

\begin{proof}
The exact Green identity gives
\(GD^{\mathsf T}D=I-\Pi_{12}\). The verified boundary complex has
\(\ker C=\operatorname{im}D\) and
\(\operatorname{im}C=\mathbf1_{20}^{\perp}\), hence
\(C^{\mathsf T}HC=I-DGD^{\mathsf T}\) and \(CC^{\mathsf T}H=I-\Pi_{20}\).
Block multiplication proves the three displayed identities.
Closedness of a cone two-cochain with boundary value \(B\) requires
\(B=Cr\) for its radial component \(r\), equivalent to zero total flux.
The radial potential \(GD^{\mathsf T}A\) has zero mean and changes by
\(\chi-\overline\chi\) under \(A\mapsto A+D\chi\), exactly matching
the apex scalar value \(\overline\chi\).

In barycentric coordinates, the local forms are
\begin{gather*}
 w_i=\lambda_i,\qquad
 w_{ij}=\lambda_i\,d\lambda_j-\lambda_j\,d\lambda_i,\\
 w_{ijk}=2\sum_{\mathrm{cyc}(i,j,k)}
                  \lambda_i\,d\lambda_j\wedge d\lambda_k.
\end{gather*}
Their exterior derivatives commute with signed incidence. Their matching
traces give the assembled identities in the weak sense
\cite{ArnoldFalkWinther2010}. Vector potentials continuous and piecewise
affine in time, with constant slab scalar potentials, give the displayed fields.
For affine \(\chi(t)\), the scalar gauge change is
\(-\partial_t\chi\); differentiation and \(d^2=0\) prove the claims.
The magnetic field is continuous in time, so its derivative acquires no
additional time-interface impulse.
\end{proof}

The radial reconstruction is global on the boundary graph: its
\(12\times30\) matrix has \(240\) nonzero entries. It is a reconstruction
rule, not a derivation of local radial transport. Its cone coefficients are
determined by the boundary data and introduce no independent radial degrees
of freedom. The finite cochain and zero-flux statements are formalized in
\path{Lean/Screen/ConeCochainBridge.lean}. That implementation uses a
dual-tree section followed by the cycle projector for the radial two-form;
uniqueness among co-closed radial solutions proves agreement with
\(C^{\mathsf T}H\) on the zero-flux domain. The Whitney field and
distributional statements are the paper argument above.

\paragraph{The metric and time defects of the actual history.}
Let \(M_1,M_2\) be the Gram matrices of \(W_1,W_2\) in the supplied
Euclidean metric, with unit constitutive coefficients. Their exact
integrals use \(\int_T\lambda_i\lambda_j=|T|(1+\delta_{ij})/20\).
These positive matrices specify geometric field energy. They are distinct
from counting inner products. Extend source covectors as
\(\widetilde\rho_n=(0,\rho_n)\) and
\(\widetilde J_n=(0_{12},J_n)\), preserving the source pairing.
Slabwise constant charge has jumps at the knots. Its compatible temporal
current functional is
\(\sum_n h\widetilde J_n\delta_{t_{n+1}}\), since
\(\rho_{n+1}-\rho_n=-hD^{\mathsf T}J_n\).
This identity uses interior test functions; an extension beyond the window
uses the initial charge before zero and the terminal charge after the last
current impulse. Ordinary constant slab current has a different pairing.

Write \(\widetilde A_n=P_1A_n\), \(\widetilde\phi_n=P_0\phi_n\),
\(\widetilde E_n=P_1E_n\), \(\widetilde B_n=P_2CA_n\), and
\(Q(b)=b^{\mathsf T}M_2b\). The exactly integrated field-plus-source action is
\begin{align*}
 S_{\mathrm{pr}}=\sum_{n=0}^{N-1}\biggl\{&
 \frac h2\widetilde E_n^{\mathsf T}M_1\widetilde E_n
 -\frac h6\left[Q(\widetilde B_n)+
 \widetilde B_n^{\mathsf T}M_2\widetilde B_{n+1}+Q(\widetilde B_{n+1})\right]\\
 &+h\widetilde J_n^{\mathsf T}\widetilde A_{n+1}
 +h\widetilde\rho_n^{\mathsf T}\widetilde\phi_n\biggr\}.
\end{align*}
Let \(S_{\mathrm{right}}\) use the same electric and source terms but
right-endpoint magnetic quadrature. Direct integration and telescoping give
\begin{equation}
 S_{\mathrm{pr}}-S_{\mathrm{right}}
 =\frac h4\left[Q(\widetilde B_N)-Q(\widetilde B_0)\right]
 +\frac h{12}\sum_{n=0}^{N-1}Q(\widetilde B_{n+1}-\widetilde B_n).
 \label{eq:cone-time-action-defect}
\end{equation}
The second term equals
\(h^2\int Q(\partial_t\widetilde B)\,dt/12\).
It is second order on a fixed time window only under uniform temporal regularity; the total action
defect includes the displayed endpoint term. Fixed-endpoint variations
remove that term. Nonuniform time grids need additional terms.
The separate spatial defect relative to the counting action is
\[
 S_{\mathrm{right}}-S_{\mathrm{count}}
 =\frac h2\sum_n\left[
 E_n^{\mathsf T}(P_1^{\mathsf T}M_1P_1-I)E_n
 -B_{n+1}^{\mathsf T}(P_2^{\mathsf T}M_2P_2-I)B_{n+1}\right].
\]
The clock action is excluded from all three field-plus-source actions.

For the derivatives, allow every cone potential coefficient to vary
independently and evaluate afterward at the reconstructed history; variations
need not stay in the images of \(P_1\) and \(P_0\).
With fixed vector-potential endpoints and free spatial traces, the full
finite-element derivatives of \(S_{\mathrm{pr}}\) are
\begin{align*}
 \frac{\partial S_{\mathrm{pr}}}{\partial\widetilde\phi_n}
 &=h\left(\widetilde\rho_n-\widetilde D^{\mathsf T}M_1\widetilde E_n\right),\\
 \frac{\partial S_{\mathrm{pr}}}{\partial\widetilde A_m}
 &=M_1(\widetilde E_m-\widetilde E_{m-1})+h\widetilde J_{m-1}\\
 &\quad-\frac h6\widetilde C^{\mathsf T}M_2
 (\widetilde B_{m-1}+4\widetilde B_m+\widetilde B_{m+1}),
 \qquad 0<m<N.
\end{align*}
The second line differs from the metric right-endpoint derivative by
\(h\widetilde C^{\mathsf T}M_2(2\widetilde B_m-
\widetilde B_{m-1}-\widetilde B_{m+1})/6\).
This temporal correction does not remove the spatial metric/source defect.
The coefficient identities and telescoping are formalized in
\path{Lean/Screen/WhitneyTimeBridge.lean}; analytic integration and
distributional source interpretation are stated here separately.

The authenticated serial history gives the following dimensionless values
for both endpoint-fixed gauge representatives:
\begin{center}
\begin{tabular}{lr}
Counting field-plus-source action & $-14.42982115034$\\
Volume action, right-endpoint magnetic term & $-14.25989506948$\\
Exactly time-integrated volume action & $-10.21278655268$\\
Spatial action defect & $0.16992608086$\\
Temporal action defect & $4.04710851680$\\
Maximum absolute scalar residual, divided by $h$ & $0.64544633271$\\
Maximum absolute interior vector residual & $1.57903507514$
\end{tabular}
\end{center}
The rational cochain identities are exact. The volume Gram matrices and
displayed values use double-precision analytic moments and independent
tetrahedral quadrature, compared with absolute and relative tolerances
\(10^{-10}\); they are not interval certificates. The independent verifier
differentiates all \(42+2\cdot13=68\) free coefficients of the action,
including the twelve radial vector variations omitted by a boundary-only
test. A separate exact counting-metric control gives temporal action defect
\(663629599/63369648\) and first vector residual \(44749/55152\).
These are mathematical controls on recorded inputs, not measurements.

Whitney one-forms have matching tangential traces; their normal jumps can
contribute to electric divergence. Magnetic tangential jumps and temporal
electric jumps likewise contribute to the sourced equations. The assembled
weak derivatives retain these effects; element-interior tests are
insufficient. The source covectors specify finite-element test functionals,
not arbitrary-test physical densities. Under the conventional smooth
potential signs, \(D^{\mathsf T}\) represents weak negative divergence, so
the coupling \(+J\cdot A+\rho_{\mathrm{load}}\phi\) corresponds to
\(\rho_{\mathrm{physical}}=-\rho_{\mathrm{load}}\) if a density
identification is made.
The supplied cone yields continuous-coordinate fields and explicit
defects. It supplies neither sourced stationary volume dynamics nor a
source-selected geometry, physical clock, spatial refinement family,
intertwiner with the translation-symbol field group, or laboratory law.
